\section{建模}

并非每个分析都需要复杂模型。从\textbf{可解释的基线}（interpretable baseline）开始。

\subsection{建模最佳实践}

对于高速公路房产问题，合理的第一步是回归或其他透明模型，在控制面积、房龄、位置等因素后，估计高速公路邻近度与房产价值之间的关系。

要求 Codex 明确以下要素：

\begin{enumerate}[nosep]
    \item \textbf{目标变量和特征定义}（target variable and feature definitions）
    \item \textbf{纳入哪些控制变量及原因}（which controls and why）
    \item \textbf{数据泄漏风险和排除项}（leakage risks and exclusions）
    \item \textbf{划分方式、评估方法和不确定性估计}（split, evaluation, uncertainty）
    \item \textbf{用通俗语言解释结果含义}（plain language interpretation）
\end{enumerate}

\begin{warningbox}{弱模型也有价值}
如果第一个模型效果不佳，这本身就是有用的信号。它能告诉你问题出在模型选择、特征工程、数据合并质量还是问题定义本身。不要因为结果不好就直接跳到复杂模型。
\end{warningbox}

\subsection{本章小结}
建模应从可解释的基线出发，要求 Codex 明确目标、特征、控制变量、泄漏风险和结果解读。弱模型是诊断问题的重要信号。

